% ===================================================================
\section{Topology Programmability}
\label{sec:programmability}

Sections~\ref{sec:tesseract} and~\ref{sec:attractor} established that
the Steersman programs polytope-derived topologies and that the
spectral attractor provides the unprogrammed baseline. This section
asks a sharper question: does the Steersman program \emph{any} pair
specification, or does the pair specification require geometric
coherence?

Two negative controls bracket the answer: wrong-labels (random
partition, WL-001) and octahedral contrastive (minimum geometry,
O-001).

\subsection{Wrong-Labels: Random Partition (WL-001)}

The wrong-labels experiment uses $n = 6$ channels with a random pair
partition instead of the RD face-pair geometry. All other parameters
(model, seed, hyperparameters, Steersman configuration) are identical
to the standard contrastive experiments.
If the Steersman is a brute-force topology programmer, the random
partition should produce block-diagonal structure aligned with the
random pairs. If geometric coherence is required, no block structure
should emerge.

\begin{table}[H]
\centering
\caption{Wrong-labels control ($n = 6$, random pair partition) vs.\
  standard contrastive ($n = 6$, RD face pairs). Both experiments use
  identical hyperparameters. The wrong-labels co/cross value reports
  alignment with the \emph{random} pair specification.}
\label{tab:wronglabels}
\small
\begin{tabular}{lccccc}
\toprule
Experiment & Pairs & Steps & Fiedler & $\rho$ & BD \\
\midrule
H-ch6 (RD) & $(0,5), (1,4), (2,3)$ & 10{,}000 & 0.00009 & 70{,}404:1 & Yes \\
WL-001 (random) & random partition & 10{,}000 & $1.26 \times 10^{-5}$ & $8.7 \times 10^{-6}$:1 & \textbf{No} \\
\bottomrule
\end{tabular}
\end{table}

No block-diagonal structure emerged at any point during 10{,}000
steps of wrong-labels training. The co-axial/cross-axial coupling
ratio reached $8.7 \times 10^{-6}$:1---effectively zero. For
comparison, the standard contrastive control (H-ch6) exceeded 200:1
by step~200 and converged to 70{,}404:1.

The Steersman drove Fiedler to $1.26 \times 10^{-5}$---below the
spectral attractor at ${\sim}0.09$---indicating that the contrastive
loss produced \emph{some} structural effect. But the effect was pure
suppression of connectivity, not block-diagonal organization. The
bridge decoupled without organizing: it lost connectivity without
gaining topology.

This is a precise failure mode. The contrastive loss drives down
both co-axial and cross-axial coupling when the pair specification
lacks geometric coherence. Without a coherent spatial decomposition,
there is no gradient signal distinguishing co-axial from cross-axial
entries---the contrastive pressure distributes uniformly and
suppresses all coupling toward zero. The result is a bridge that
is neither connected (Fiedler $\approx 0$) nor structured
(co/cross $\approx 0$)---worse than the unprogrammed attractor.

\subsection{Octahedral Contrastive: Minimum Geometry (O-001)}

The octahedral experiment tests the other boundary: the simplest
possible geometric specification. With $n = 4$ and only 2~co-axial
pairs against 4~cross-axial pairs, the octahedral topology provides
the minimum spatial structure that could plausibly produce
block-diagonal organization.

\begin{table}[H]
\centering
\caption{Octahedral contrastive ($n = 4$, 2 co-axial pairs from
  octahedral vertex opposition) at 10{,}000 steps. Per-bridge mean
  ratio reported; pooled ratio is 369{,}365:1.}
\label{tab:octahedral}
\small
\begin{tabular}{lccccc}
\toprule
Experiment & $n$ & Pairs & Fiedler & $\rho$ & BD \\
\midrule
H-ch4 (spectral-only) & 4 & none & 0.0836 & ${\sim}1$:1 & No \\
O-001 (octahedral) & 4 & $(0,1), (2,3)$ & $1.1 \times 10^{-5}$ & \textbf{473{,}622:1} & \textbf{Yes} \\
\bottomrule
\end{tabular}
\end{table}

O-001 produces the strongest block-diagonal signal in the entire
research programme (Figure~\ref{fig:o001_emergence}): per-bridge mean co-axial/cross-axial coupling
ratio of \textbf{473{,}622:1} (pooled: 369{,}365:1, median:
401{,}851:1, range: 126{,}782--1{,}577{,}518:1 across 88~bridges).
This exceeds both the RD result (70{,}404:1 at $n = 6$) and the
tesseract result (41{,}564:1 at $n = 8$) by an order of magnitude.

\begin{figure}[t]
  \centering
  \includegraphics[width=\textwidth]{figures/fig_o001_emergence.pdf}
  \caption{Octahedral block-diagonal emergence (O-001, $n = 4$).
    Co-axial/cross-axial coupling ratio grows monotonically to
    473{,}622:1---the programme's strongest signal. The minimum
    geometry (two co-axial pairs, four cross-axial pairs) produces
    maximum coupling strength, establishing the inverse-$n$
    relationship.}
  \label{fig:o001_emergence}
\end{figure}

The inverse relationship between channel count and coupling ratio
is established across three data points: $n = 4$ (${\sim}474$K:1)
$>$ $n = 6$ (${\sim}70$K:1) $>$ $n = 8$ (${\sim}42$K:1). Fewer channels means fewer
cross-axial entries to suppress: $n = 4$ has only 4~cross-axial
pairs, while $n = 8$ has 24. The co-axial signal is concentrated
into fewer entries, producing stronger per-pair coupling and a
more extreme ratio. The co-planar coupling magnitude at O-001 also
grows linearly throughout training ($0.086$ at step~500, $1.027$ at
step~10{,}000) with no saturation, indicating the ratio would
increase further with continued training.

\subsection{The Geometric Coherence Requirement}

The wrong-labels and octahedral results, taken together, define the
boundary of topology programmability. On one side: any polytope-derived
spatial specification---RD at $n = 6$, tesseract at $n = 8$, octahedron
at $n = 4$---produces block-diagonal structure with coupling ratios
from 42{,}000:1 to 200{,}000:1. On the other side: a random partition
at the same channel count ($n = 6$) produces no structure at all.

The Steersman is not a brute-force topology enforcer. It is a
\emph{selective geometric compiler}. It compiles valid geometric
programs---pair specifications that encode the co-axial face or vertex
structure of a polytope---and rejects incoherent specifications by
driving the bridge toward total decoupling. The pair specification is
not a suggestion; it is a program that must be well-formed.

This selectivity is a feature, not a limitation. It means the
block-diagonal structure observed under polytope-derived specifications
is genuine geometric encoding, not an artifact of the contrastive loss
formulation. The loss function is identical in WL-001 and H-ch6; only
the pair indices differ. The indices carry the geometric information,
and the Steersman distinguishes coherent from incoherent geometry.

\subsection{Block-Diagonal Hierarchy}

Table~\ref{tab:hierarchy} presents the complete block-diagonal hierarchy
across all experiments, from the strongest positive to the strongest
negative.

\begin{table}[H]
\centering
\caption{Complete topology hierarchy across all experiments, ordered
  by co-axial/cross-axial coupling ratio. Four regimes emerge,
  separated by orders of magnitude.}
\label{tab:hierarchy}
\small
\begin{tabular}{llcccl}
\toprule
Experiment & Topology & $n$ & $\rho$ & Fiedler & Regime \\
\midrule
O-001 & Octahedral & 4 & 473{,}622:1 & $1.1 \times 10^{-5}$ & Geometric BD \\
Seed-43 & RD contrastive & 6 & 73{,}309:1 & ${\sim}10^{-5}$ & Geometric BD \\
H-ch6 & RD contrastive & 6 & 70{,}404:1 & $9 \times 10^{-5}$ & Geometric BD \\
T-001r2 & Tesseract & 8 & 41{,}564:1 & $1.9 \times 10^{-4}$ & Geometric BD \\
\midrule
E-001 & Emanation & 6 & 1.12:1 & 0.084 & Hierarchical \\
\midrule
H-ch3/4/8/12 & Spectral-only & 3--12 & ${\sim}1$:1 & 0.084--0.102 & Spectral attractor \\
\midrule
WL-001 & Wrong-labels & 6 & $8.7 \times 10^{-6}$:1 & $1.3 \times 10^{-5}$ & Collapse \\
R-001 & Resonance & 6 & $8.9 \times 10^{-6}$:1 & $1.2 \times 10^{-5}$ & Collapse \\
\bottomrule
\end{tabular}
\end{table}

The hierarchy separates into four regimes with no overlap
(Figure~\ref{fig:four_regimes}).

\begin{figure}[t]
  \centering
  \includegraphics[width=\textwidth]{figures/fig_four_regimes.pdf}
  \caption{Four dynamical regimes of the bridge matrix under different
    programming conditions, ordered by co-axial/cross-axial coupling
    ratio. Block-diagonal ($\rho > 40{,}000$:1), hierarchical coherence
    ($\rho \approx 1$:1, Fiedler $\approx 0.08$), spectral attractor
    ($\rho \approx 1$:1, Fiedler $\approx 0.09$), and connectivity
    collapse ($\rho \approx 10^{-5}$:1). No intermediate outcomes
    exist between regimes.}
  \label{fig:four_regimes}
\end{figure}

\textbf{Regime 1: Geometric block-diagonal} ($\rho > 40{,}000$:1).
Polytope-derived pair specifications (octahedral, RD, tesseract)
produce extreme co-axial/cross-axial separation with near-zero
Fiedler. The regime spans one order of magnitude internally
(42K--474K) and is separated from the next regime by four orders
of magnitude.

\textbf{Regime 2: Hierarchical coherence} ($\rho \approx 1$:1,
Fiedler $\approx 0.08$). The emanation architecture (E-001) produces
slight directional preference without block-diagonal structure. The
bridge maintains connectivity near the spectral attractor but
develops weak asymmetry through the master-offset hierarchy
(Section~\ref{sec:nongeometric}).

\textbf{Regime 3: Spectral attractor} ($\rho \approx 1$:1,
Fiedler $\approx 0.09$). Spectral-only training (no contrastive
loss) produces connected, isotropic bridges. This is the
unprogrammed state (Section~\ref{sec:attractor}).

\textbf{Regime 4: Connectivity collapse} ($\rho \approx 0$:1,
Fiedler $\approx 0$). Incoherent contrastive specifications
(wrong-labels, resonance) produce total collapse: both
co-axial/cross-axial ratio and Fiedler approach zero. The bridge
disconnects without organizing. This regime is distinct from the
spectral attractor: the spectral attractor maintains stable
connectivity ($\text{Fiedler} \approx 0.09$), while the collapse
regime drives connectivity to zero ($\text{Fiedler} \approx 10^{-5}$).
The contrastive loss, given incoherent pairs, overwhelms the spectral
loss and suppresses all coupling.

The four regimes are cleanly separated. No experiment falls between
regimes. The Steersman's response to a pair specification is
categorical: geometric coherence produces block-diagonal structure;
hierarchical indirection produces weak coherence; no specification
produces the spectral attractor; incoherent specification produces
collapse.
